\documentclass{rep}

\title{Proyecto 2}
\author{Francisco Uriel Pineda Hérnandez}

\date{\today}

\subtitle={Reporte técnico de aplicación: Base de Datos Musical}
\email={\href{mailto:ciscohdz@ciencias.unam.mx}{\color{principaluno}\texttt{ciscohdz@ciencias.unam.mx}}}
\subject={Modelado y programación}

\setlength{\columnsep}{14pt}

% --- Archivo de bibliografía ---
\addbibresource{references.bib}

% --- Inicio del documento ---
\begin{document}
	
	\pagestyle{fancy}
	\unspacedoperators
	
% --- Título ---
	\twocolumn[
		\begin{center}
			\maketitle
			
			{\begin{tcolorbox}[colframe=white, colback=principaldos, arc=8pt]
				\begin{onecolabstract}

                  La gestión de colecciones locales de música digital
                  suele verse limitada por metadatos desorganizados y
                  por herramientas que alteran los ficheros de audio
                  originales al modificar su información. Para
                  resolver esta problemática, se desarrolló una
                  aplicación orientada a la catalogación, búsqueda y
                  edición de bibliotecas musicales en formato MP3
                  mediante el uso de una base de datos relacional. En
                  el sistema se implementa un 'minero' de directorios
                  capaz de extraer etiquetas ID3v2.4, una capa de
                  persistencia con operaciones para inserción,
                  actualización y busqueda de registros sin necesidad
                  de escribir sentencias SQL, así como una interfaz
                  gráfica que permite editar metadatos y clasificar
                  intérpretes entre individuos y agrupaciones
                  musicales. La solución fue construida bajo el
                  sistema de construcción Meson utilizando los
                  lenguajes Vala y C apoyándose en las bibliotecas
                  GLib, GObject y GIO para el modelo de objetos y
                  eventos; GTK4 para la interfaz de usuario; libid3tag
                  para la lectura de metadatos; y SQLite3 como motor
                  de almacenamiento. Como resultado, se obtuvo un
                  gestor musical que sincroniza en tiempo real las
                  modificaciones realizadas en la base de datos con
                  las vistas gráficas.
                  
				\end{onecolabstract}
			\end{tcolorbox}}

			\medskip
		\end{center}
	]

	\selectlanguage{spanish}
	
% --- Cuerpo del reporte ---
	
	\section{Introducción}

    El objetivo de este segundo proyecto es crear un sistema capaz de
    crear y manejar una base de datos musical, se especifica el uso
    de la biblioteca SQLite3 como motor para la base de datos.

    A través de una interfaz gráfica, el sistema podrá recibir un
    diectorio que almacene archivos musicales con etiquetas id3v2.4 y
    que, al recorrerlo, obtenga la información con la que poblar
    nuestra base de datos.Además, se busca que el usuario pueda
    explorar su colección, actualizar registros y clasificar
    intérpretes sin alterar los archivos MP3 en disco y sin requerir
    conocimientos del lenguaje SQL para efectuar búsquedas.
    
    \section{Modelado}

    Habiendo establecido las necesidades del sistema y los objetivos a
    cumplir, modelaremos la solución analizando detalladamente los
    requerimientos para dividir cada problema en pequeños fragmentos
    más manejables para, posteriormente, diseñar componentes que los
    resuelvan.

    \subsection{Análisis}

    De manera general, la primera división surge de la necesidad de
    aislar las responsabilidades del sistema. Esto nos permite
    garantizar que las operaciones de entrada, procesamiento y salida
    estén repartidas y delitmitadas. 
    
    El sistema tiene que resolver: el minado (como denominamos a la
    obtención de datos de archivos mp3); la edición de los registros
    en la base datos; la búsqueda de registros en la base de datos; la
    reproducción de contenido multimedia; el manejo de las operaciones
    sobre la base de datos; asi como la visualización y entrada de los
    datos provenientes del usuario.

    De modo que dichas responsabilidades se pueden abstraer en los
    componentes: Vista, Minero, Editor, Buscador, Reproductor y
    Manejador de la base de datos.

    Ahora detallaremos con mayor profundidad cada uno de los
    componentes para entender los problemas que buscan resolver.

    \subsubsection{Vista}
    
    La \textbf{Vista}, al ser la parte más cercana al usuario, debe permitir
    interacciones tanto de entrada como de salida.

    Se deberá poder permitir:
    \begin{enumerate}
    \item Visualizar y editar la información de las canciones y álbumes
    \item Definir a un intérprete como persona o grupo
    \item Agregar personas a grupos
    \item Consultar títulos, intérpretes o álbumes de las canciones.
    \end{enumerate}

    En ese sentido debe proveer información clara sobre el estado
    actual de la base de datos. Además los elementos interactivos
    tiene que estar organizados de manera que se pueda diferenciar las
    funcionalidades que provee el sistema.
    
    \subsubsection{Minero}

    Como se ha especificado, el \textbf{Minero} tiene que recorrer un
    directorio; por cada archivo mp3 con etiquetas id3v2.4 que
    encuentre debe obtener los metadatos de dicha canción y enviarlos
    al Manejador de la base de datos para que sean ingresado a la misma.
    
    Se intentará leer las siguientes etiquetas ID3v2.4:

    \begin{enumerate}
    \item \textbf{TPE1}, \textit{Lead performer(s)/Soloist(s)}, el
      intérprete de la canción.
    \item \textbf{TIT2}, \textit{Title/songname/content} description,
      el título de la canción.
    \item \textbf{TALB}, \textit{Album/Movie/Show title}, el nombre
      del álbum.
    \item \textbf{TDRC}, \textit{Recording time}, la fecha de
      grabación; usualmente el año basta.
    \item \textbf{TCON}, \textit{Content type}, el tipo del contenido;
      el género, normalmente.
    \item \textbf{TRCK}, \textit{Track number/Position in set}, el
      número de pista (un entero) o posición en el álbum (en el
      formato “3/10”, tercera de diez pistas).
    \end{enumerate}

    Además de las etiquetas id3, el proceso de minado también debe
    obtener datos como:
    
    \begin{enumerate}
    \item Ruta al archivo mp3
    \item Ruta al directorio que contiene al archivo mp3.
    \item Nombre del archivo mp3
    \item Año de creación del archivo mp3
    \end{enumerate}

    Cada que se agregue correctamente un registro a la base de datos
    el minero debe poder informar a la Vista los cambios que se
    hicieron en la base de datos.

    Para garantizar la consistencia entre la información ingresada a
    la base de datos con la que se presenta al usuario, el proceso de
    minado requiere un mecanismo de comunicación. Cada vez que un
    registro se inserta con éxito, el sistema debe ser capaz de emitir
    un aviso para que la Vista refleje las modificaciones. 
    
    \subsubsection{Editor}

    El \textbf{Editor} debe actualizar los registros en la base de
    datos, procesando las modificaciones que el usuario solicite a
    través de las opciones que establece la Vista.
    
    De manera general notamos que la edición de la base de datos se
    puede dividir en tres grupos: la edicion de canciones, la edición
    de álbumes y la edición de intérpretes. Editar cada una de la
    partes implica diferentes procesos, desde la forma en que se
    presenta la información a editar hasta las comprobaciones que
    deben hacerse para efectuar el cambio que se pretende hacer. 

    Si queremos editar una canción lo primero que notamos es que
    podemos modificar su titulo, \textit{track}, año de lanzamiento y
    género al que pertenece, pero también debe existir la posibilidad
    de cambiar únicamente esa canción de album así como de intérprete.

    Por su parte. al editar un álbum podemos modificar su nombre y el
    año de lanzamiento del mismo. Hacer modificaciones directamente a
    un álbum implica que todas las canciones que estén vinculadas a
    éste reflejen los cambios realizados.

    Ahora bien, para modificar un intérprete necesitamos saber si se
    trata de una persona o de un grupo. Sabiendo lo anterior, si el
    intérprete resulta ser una persona podemos editar el nombre
    artístico, su nombre real, su fecha de nacimiento o de defunción
    (si aplica); si es un grupo podemos editar su nombre, su fecha de
    inicio o de término (si aplica). Hacemos notar que los cambios de
    persona a grupo y viceversa deben ser posibles. 

    Una vez el usuario determine qué cambios hacer y con qué datos, el
    Editor solicita al Manejador que se apliquen directamente en la
    base de datos e informa a la Vista sobre aquellos que sí fueron
    realizados para su respectiva actualización.
    
    \subsubsection{Buscador}

    La responsabilidad del \textbf{Buscador} radica en interpretar
    parámetros de consulta provistos desde la interfaz gráfica y
    traducirlos en peticiones de filtrado que el Manejador pueda
    procesar.

    Dadas las relacionales en la información musical, una consulta no
    se limita únicamente a las propiedades directas de una
    canción. Realizar búsquedas permite al usuario utilizar criterios
    combinados y dependencias (como el nombre del álbum, el género o
    la clasificación del intérprete) para refinar los resultados
    obtenidos.

    \subsubsection{Manejador de la Base de Datos}

    El \textbf{Manejador} debe servir como el único punto de contacto
    con la base de datos, proporcionando a los demás componentes una
    estandarizada para insertar, actualizar y consultar información
    sin que estos necesiten conocer los detalles. 

    \subsection{Diseño}

    Habiendo establecido las responsabilidades de cada módulo en la
    etapa de análisis, es necesario buscar una manera de organizar la
    interacción entre los mismos. Para resolver este problema, la
    estructura de la aplicación se fundamenta en el patrón de diseño
    Modelo-Vista-Controlador (MVC).

    La adopción del patrón MVC se justifica bajo el principio de que
    la información gestionada por el programa debe estar aislada de su
    representación y de los mecanismos de interacción del usuario
    \cite{becerril}. De acuerdo con este paradigma, el sistema se
    divide en tres capas fundamentales: el \textit{modelo}, encargado
    de encapsular los datos puros y los algoritmos con que se operan;
    la \textit{vista}, que constituye la salida inmediata y perceptiva
    hacia el usuario; y el \textit{controlador}, que actúa como
    intermediario entre los otros dos. Este último gestiona las
    entradas de la vista para actualizar el modelo y se asegura de que
    cualquier cambio asíncrono en los datos se refleje
    gráficamente. Esta separación evita que el modelo se contamine con
    lógicas de interfaz y facilitando su portabilidad \cite{becerril}.

     \begin{figure}[htbp]
		\centering
        \includesvg[width=0.5\textwidth]{img/div-mvc-musica.svg}
		\caption{Diagrama de la división del sistema.}
		\label{fig:div_mvc}
	\end{figure}

    Veamos como, en el contexto de nuestra base de datos musical, los
    componentes analizados previamente resuelven sus interacciones
    acoplándose a este patrón.

    \subsubsection{Modelo}

    Como se estableció previamente, el Modelo encapsula la gestión y
    el estado de la colección musical. Dada la complejidad de la
    información que debe manejar, nos apoyaremos en una serie de
    patrones de diseño que garantizan la cohesión, el bajo
    acoplamiento y la integridad de los datos.

    Para el tránsito de la información entre la base de datos, se
    implementó el patrón de diseño \textit{Data Transfer Object}
    (DTO). Este patrón se materializa como se muestra en la Figura
    \ref{fig:DTOs} mediante las clases \texttt{Rola}, \texttt{Album},
    \texttt{Performer}, \texttt{Person} y \texttt{Group} . Su diseño
    se simplifica dado que su único propósito es empaquetar conjuntos
    de atributos (como el título, el año o la ruta del archivo) para
    que puedan ser transportados en una sola llamada.

    \begin{figure}[htbp]
      \centering
        \includesvg[width=\columnwidth, inkscapelatex=false]{img/DTOs.svg}
        \caption{Diagrama UML de clases DTO.}
        \label{fig:DTOs}
    \end{figure}
    
    Dentro de esta estructura, destaca la jerarquía de herencia
    diseñada para los intérpretes musicales. Una pista de audio
    siempre es ejecutada por un intérprete genérico. No obstante,
    puede ser un individuo o un grupo. Para resolver esto, se abstrae
    la clase base \texttt{Performer} (que contiene atributos comunes
    como el identificador y el nombre general), de la cual heredan las
    clases \texttt{Person} (que añade metadatos específicos como fecha
    de nacimiento/defunción y nombre real/artístico) y \texttt{Group}
    (que incluye fecha de incio/termino y una lista de sus
    integrantes). Esta jerarquización permite aplicar polimorfismo:
    las demás entidades del sistema, como la clase \texttt{Rola},
    pueden referenciar a un \texttt{Performer} de manera genérica,
    mientras que el sistema conserva la capacidad de tratar los
    detalles específicos de personas o grupos según sea necesario.

    Ahora bien, las responsabilidades que en el análisis recaían sobre
    el \textbf{Manejador de la Base de Datos}, se delegan a la clase
    \texttt{Dao}. Esta decisión obedece al patrón \textit{Data Access
      Object} (DAO), el cual dicta que toda la lógica de acceso a
    datos (consultas, inserciones, actualizaciones y eliminaciones)
    debe aislarse de los demás componentes del sistema.

    Así, el \texttt{Dao} traduce los objetos DTO hacia la base de
    datos (y viceversa). Esto presenta una ventaja fundamental:
    componentes funcionales como el \texttt{Minero} o el
    \texttt{Editor} no necesitan conocer los detalles de la base de
    datos para funcionar. Si en el futuro se requiriera cambiar el
    motor de almacenamiento, los cambios se limitarían exclusivamente
    a la implementación de la clase \texttt{Dao}, dejando el resto del
    modelo intacto.

    Las clases \texttt{Minero} y \texttt{Editor} intervienen como se
    muestra en la Figura \ref{fig:Modelo}. La primera se encarga de
    recorrer las estructuras de directorios, extraer metadatos puros
    de los archivos multimedia y transformarlos en objetos DTO
    validos. Por su parte, el segundo centraliza la modificación de
    registros. Por ejemplo, al editar un intérprete, el Edito revalúa
    escenarios complejos: si un intérprete cambia de ser una
    \texttt{Person} a \texttt{Group}, el Editor se encarga de la
    eliminación de datos obsoletos y de la creación de nuevas
    asociaciones, delegando posteriormente las operaciones de
    persistencia en la base de datos al \texttt{Dao}.

    Para mantener separado el Modelo y evitar que este deba invocar
    directamente a la Vista o al Controlador cuando ocurre un cambio
    en los datos, se integró el patrón de diseño \textit{Observer}
    mediante la figura de un \textbf{Notificador} que pertenecerá al
    \textbf{Controlador}.

    \begin{figure}[htbp]
      \centering
        \includesvg[width=\columnwidth, inkscapelatex=false]{img/Modelo.svg}
        \caption{Diagrama UML de clases principales del Modelo.}
        \label{fig:Modelo}
    \end{figure}

    \subsubsection{Vista}
    
    La capa de presentación tiene la responsabilidad de traducir la
    información provista por el Modelo en elementos visuales, así como
    de capturar las interacciones del usuario. La interfaz gráfica se
    diseñó dividiendo el espacio de trabajo en dos secciones
    funcionales: superior e inferior.

    La barra superior determina el flujo principal de la
    aplicación. Se subdivide lógicamente en tres segmentos:

    \begin{itemize}
    \item \textbf{Izquierda:} Aloja los controles de navegación
      principal, permitiendo alternar entre vistas rápidas: la vista
      general de Canciones, la colección de Álbumes y el catálogo de
      Intérpretes. Además, contiene las acciones de entrada,
      destacando el botón para accionar el proceso de minado (o
      re-minado) de un directorio, y el acceso a la herramienta de
      consulta (búsqueda).
    \item \textbf{Centro:} Reservado para la identidad de la
      aplicación, mostrando el título del sistema.
    \item \textbf{Derecha:} Incluye el control para finalizar la
      ejecución del programa.
    \end{itemize}

    La parte inferior de la interfaz es la encargada de mostrar los
    resultados de las consultas o el contenido de las vistas
    rápidas. Para optimizar la experiencia de usuario, la disposición
    de los datos cambia dependiendo de la entidad que se esté
    visualizando:
    
    \begin{itemize}
    \item \textbf{Vista de Canciones (Principal):} Muestra el
      contenido en un formato de lista secuencial, donde cada fila
      expone los metadatos esenciales en el orden: imagen genérica,
      título, intérprete, álbum y género.
    \item \textbf{Vista de Álbumes:} Cambia a una disposición en
      cuadrícula, exhibiendo tarjetas que contienen una imagen
      cuadrada representativa, el nombre del álbum y su respectivo
      año.
    \item \textbf{Vista de Intérpretes:} Mantiene la cuadrícula
      anterior, pero presenta un formato de imagen circular acompañado
      únicamente por el nombre del intérprete.
    \end{itemize}

    La manera en que se proveen los resultados en el contenedor
    central está íntimamente ligada a las capacidades de consulta. La
    barra de búsqueda actúa en las tres vistas. Mientras que en la
    vista de canciones se aprovecha todos los campos que identifican a
    las mismas, en las vistas de álbumes e intérpretes el filtro opera
    sobre los rubros expuestos.

    Para cumplir con los requerimientos de actualización de registros,
    se evitó saturar la interfaz principal con formularios
    estáticos. En su lugar, se optó por un diseño basado en
    interacciones contextuales. Por comodidad y agilidad, el usuario
    puede realizar un clic derecho sobre cualquier elemento gráfico
    (una fila de canción, una tarjeta de álbum o de intérprete) para
    desplegar un menú emergente. Al seleccionar la opción de
    ``Editar'', el sistema lanza un diálogo flotante que contiene los
    campos específicos de la entidad seleccionada, permitiendo
    modificar sus atributos de forma aislada antes de solicitar la
    actualización al Editor.

    \begin{figure}[htbp]
      \centering
        \includesvg[width=\columnwidth, inkscapelatex=false]{img/ClasesVista.svg}
        \caption{Diagrama UML de clases principales de la Vista.}
        \label{fig:Vista}
    \end{figure}
    
    \subsubsection{Controlador}

    El propósito fundamental del \textbf{Controlador} es garantizar
    una estricta separación entre el Modelo y la Vista.  Para lograr
    este nivel de desacoplamiento, el diseño del controlador se divide
    funcionalmente en dos mecanismos principales: la gestión de
    interacciones de entrada y la propagación de eventos de salida.

    El Controlador asume la responsabilidad de interceptar los eventos
    generados por el usuario desde la Vista, tales como la petición de
    minar un directorio o la solicitud de edición sobre canciones,
    álbumes e intérpretes. Al capturar estas acciones, el sistema
    delega la ejecución de la lógica al Modelo, conectando los eventos
    de la interfaz con las operaciones de actualización
    correspondientes.

    Por otro lado, para notificar a la Vista sobre los cambios
    ocurridos en el Modelo sin que este último deba conocer la
    existencia de la interfaz, se consolidó el patrón
    \textit{Observer} a través de un componente centralizado de
    notificación al cual llamaremos \textbf{Notificador}.
    
    Este mecanismo establece el uso de orientación a eventos. Cuando
    el \texttt{Minero} inserta exitosamente una nueva canción, o
    cuando el \texttt{Editor} finaliza la actualización de un objeto,
    estos le indican al Notificador que emita una señal. Cualquier
    componente externo interesado en el estado de la base de datos
    puede suscribirse a estas notificaciones.

    Dada la naturaleza de la información musical, las operaciones
    suelen tener efectos secundarios. Por ejemplo, al insertar una
    nueva canción o editar los metadatos de una existente, es posible
    que implícitamente se descubra un nuevo álbum o intérprete. Para
    manejar estos escenarios, el diseño integra estructuras de
    resultados encapsuladas que viajan junto con las
    notificaciones. Estos objetos de estado empaquetan la entidad
    procesada junto con banderas que indican, con precisión, si la
    operación originó la creación de registros adicionales.
    
    \section{Implementación}

    En esta sección hablaré de las cosas que, para mí, fueron más
    complejas o que más me sorprendieron. Para ello, detallaremos las
    bibliotecas y mecanismos del lenguaje empleados, así como la forma
    en que se resolvieron algunos de los problemas técnicos propios de
    integrar Vala, C, GTK4 y SQLite3.

    \subsection{Elección de tecnologías}

    Para empezar, se optó por utilizar \textbf{Vala} como lenguaje
    principal de desarrollo en conjunto con \textbf{C}. Lo anterior ya
    que Vala permite programar bajo un paradigma orientado a objetos,
    compilando directamente a código C sobre el sistema de tipos
    \texttt{GObject}, mientras que no queria dejar pasar la
    oportunidad de utilizar archivos VAPI (\textit{Vala Application
      Programming Interface}) para vincular código C directamente en
    Vala.
    
    Nuevamente se usa \textbf{Meson} como sistema de construcción.
    Nos permitió abilitar los compiladores de C y Vala, declarar las
    dependencias \texttt{glib-2.0}, \texttt{gobject-2.0} y
    \texttt{gio-2.0} para las estructuras de datos, hilos y manejo de
    archivos; \texttt{gtk4} para la interfaz gráfica; \texttt{sqlite3}
    como motor de base de datos; y \texttt{id3tag}
    (\texttt{libid3tag}) para la lectura de metadatos de archivos de
    mp3.

    Adicionalmente, para evitar que el programa dependa de rutas
    relativas externas al ejecutarse desde cualquier directorio, se
    empleó el sistema de empaquetado de recursos
    \texttt{GResource}. Durante la compilación, tanto las plantillas
    XML de la interfaz gráfica como el archivo con el esquema inicial
    de la base de datos (\texttt{rolas.sql}) se compilan e incrustan
    dentro del binario final, permitiendo cargarlos directamente desde
    memoria en tiempo de ejecución (Una maravilla).

    \subsection{Manejo de base de datos (SQLite3)}

    Toda la interacción con la base de datos se centralizó en la clase
    \texttt{Dao}, implementada bajo el patrón \textit{Singleton} para
    garantizar un único punto de acceso a la conexión.
    
    Todas las operaciones de inserción, consulta, actualización y
    eliminación se implementaron utilizando sentencias preparadas
    (\texttt{Sqlite.Statement}), lo que nos permite enlazar cada
    parámetro de la sentencia SQL de manera arbitraria, estandarizando
    la forma en que se realizan las consultas. Esto asegura la
    integridad de los tipos de datos a consultar.

    Uno de los retos técnicos más relevantes consistió en mapear
    intérpretes (bien como persona o grupo) hacia tablas sin perder
    información. Para resolver esto, las consultas de lectura como
    \texttt{buscar\_rolas()} y \texttt{buscar\_performers()} emplean
    cláusulas \texttt{LEFT JOIN} entre la tabla base
    \texttt{performers} y las tablas especializadas \texttt{persons} y
    \texttt{groups}. Al recorrer cada fila resultante, el sistema
    inspecciona el atributo \texttt{id\_type} e instancia la subclase
    concreta apoyándose en el patrón \textit{Builder}
    (\texttt{PersonBuilder}, \texttt{GroupBuilder} y
    \texttt{PerformerBuilder}).

    Asimismo, para permitir la existencia de personas que no
    interpretan canciones y la pertenencia de integrantes en grupos,
    se gestionaron los vínculos a través de la tabla intermedia
    \texttt{in\_group}. Cuando el usuario confirma una edición, la
    clase \texttt{Editor} coordina con el \texttt{Dao} la inserción o
    actualización de cada entidad, así como la sincronización de
    miembros (\texttt{agregar\_person\_in\_group()} y
    \texttt{eliminar\_person\_de\_group()}) sin eliminar los registros
    base de la tabla \texttt{persons}.

    \subsection{Minería de archivos MP3}

    La extracción de metadatos fue implementada en el componente
    \texttt{Minero} apoyándose en las utilidades de sistema de
    archivos de \texttt{GIO} (\texttt{GLib.File} y
    \texttt{FileEnumerator}) y en la biblioteca \texttt{libid3tag}. Al
    recibir un directorio, el minero recorre iterativamente su
    estructura. Por cada archivo válido, infiere automáticamente dos
    rutas fundamentales: la ruta absoluta hacia el archivo de audio
    para la entidad \texttt{Rola}, y la ruta del directorio contenedor
    para asociarla a la entidad \texttt{Album}.

    Para leer las etiquetas ID3v2.4 sin modificar los archivos
    originales, cada pista se abre en modo de solo lectura
    (\texttt{ID3\_FILE\_MODE\_READONLY}). A partir de la etiqueta
    obtenida, se extraen los marcos ya especificados en el
    Análisis. Antes de insertar una nueva canción, el minero consulta
    al \texttt{Dao} para reutilizar álbumes o intérpretes ya
    existentes y evitar duplicados.

    El controlador orquesta en un hilo de ejecución diferente el
    proceso de minado. De este modo, se evita que las operaciones de
    lectura en disco y escritura en la base de datos congelen la
    interfaz gráfica. Dado que GTK4 no permite modificar widgets
    directamente desde hilos secundarios, cada vez que el minero
    inserta un registro, la clase \texttt{Notificador} programa la
    emisión de la señal en el hilo principal de eventos utilizando
    \texttt{GLib.Idle.add()}.

    \subsection{Interfaz gráfica en GTK4}

    El diseño final de la interfaz se implementó con el software de
    creación de interfaces de usuario: Cambalache. Usando las
    anotaciones de Vala \texttt{[GtkTemplate]} y \texttt{[GtkChild]},
    se nos permite referenciar directamente los componentes visuales
    definidos en los archivos \texttt{.ui} (generados por Cambalache)
    sin saturar el código fuente con instanciación manual de
    contenedores.

    En la ventana principal, la visualización de datos se resolvió
    mediante un \texttt{Gtk.ListBox} (\texttt{lista\_rolas}) para el
    listado secuencial de canciones y dos contenedores
    \texttt{Gtk.FlowBox} para las cuadrículas de álbumes e
    intérpretes.

    Cuando el usuario escribe una consulta, el sistema construye
    contenedores temporales con los resultados filtrados y reemplaza
    dinámicamente la propiedad \texttt{child} de los
    \texttt{Gtk.ScrolledWindow} correspondientes; al limpiar la barra
    de búsqueda, se restituyen las referencias originales de manera
    inmediata.

    Las acciones de edición se integraron mediante controladores de
    eventos \texttt{Gtk.GestureClick} configurados para responder al
    botón derecho del ratón sobre cada fila o tarjeta desplegando un
    menú contextual \texttt{Gtk.Popover} posicionado en las
    coordenadas del clic.

    Al seleccionar la edición de un intérprete, se despliega
    \texttt{DialogoEdicionPerformer}, el cual enlaza un selector
    \texttt{Gtk.DropDown} con un contenedor \texttt{Gtk.Stack} para
    alternar dinámicamente entre los formularios
    \texttt{vista\_person}, \texttt{vista\_group} y
    \texttt{vista\_unknown}.

    En el caso de los grupos, el diálogo mantiene una lista temporal
    en memoria (\texttt{integrantes\_tmp}) sincornizada con un
    \texttt{Gtk.ListBox} donde cada fila incorpora un botón para
    eliminar integrantes antes de confirmar la operación. 

    Finalmente, la comunicación entre el Modelo y la Vista se
    materializó a través de la clase \texttt{Notificador}, la cual
    utiliza señales nativas de \texttt{GLib.Object}
    (\texttt{rola\_agregada}, \texttt{rola\_editada},
    \texttt{album\_editado} y \texttt{performer\_editado}). Junto con
    estas señales viajan objetos como \texttt{ResultadoInsertar} y
    \texttt{ResultadoEditarRola}, lo que permite a \texttt{Ventana}
    saber con exactitud si una edición o inserción generó un nuevo
    álbum o intérprete y actualizar únicamente los elementos visuales
    afectados.
    
    \section{Conclusiones}

    Este segundo proyecto fue, en comparación, más díficil que el
    primero. El uso de una interfaz gráfica tan interconectada como la
    que se solicitaba aumenta el número de casos que hay que tener en
    cuenta. Por su parte el manejo de una base de datos relacional,
    que si bien cuenta con pocas operaciones básicas, se vuelve
    bastante confuso cuando las tablas aumentan en complejidad.

    Un detalle del diseño que después entorpeció mucho el desarrollo,
    fue el no separar en otras vistas las personas y los grupos. Hizo
    que la complejidad de darle uso a la tabla \texttt{in\_group}
    creciera mucho más de lo que en un principio esperaba.

    El manejo de señales en Vala fue a la vez un milagro por
    lo convenientes que resultan al manejar eventos de la interfaz,
    pero un martirio a la hora de comprender como conectarlas y
    enviarlas para realizar diversas acciones.

    Varias cosas se quedaron en ideas o en diseño preliminar. La que
    más pesa es el pequeño compilador de lenguaje natural a setencias
    SQL, a pesar de no haber podido implementarlo encontré varios
    recursos que me pueden ayudar a hacerlo en un futuro. La
    reproducción de canciones es otra de esas cosas que tuve que
    sacrificar por falta de planeación. Por último (y de la que tenía
    muchas ganas) era de obtener y mostrar las imágenes que podían
    venir como etiquetas ID3v2.4, prueba de ello es que reservé
    espacios en todo el diseño de la interfaz con la esperanza de
    poblarlo.

    Lo que más rescato en está ocasión es lo mucho que pude investigar
    sobre aplicaciones más enfocadas la experiencia del usuario. Me
    sorprendieron mucho cosas como el uso que se les puede dar a las
    funciones \textit{lamda}, como por ejemplo que ``guarden'' las
    referencias de objetos que luego, sin importar cuando se ejecuten
    dichas funciones, el objeto puede está dentro de ese contexto.

    Me llevo una mejor experiencia al desarrollar este proyecto, en
    parte por todo lo aprendido del pasado, pero sobre todo por el
    tipo de cosas que me permitió aprender sobre bases de datos,
    interfaces de usuario y el aprendizaje de un nuevo lenguaje de
    programación. 

    \nocite{*}
    \printbibliography[heading=bibintoc]
	
\end{document}
